%% Shared preamble of the main paper and the appendix.
\usepackage{algorithm}
\usepackage{algpseudocode}
\usepackage{multirow}
\usepackage{graphicx}
\usepackage{enumitem}
\usepackage{needspace}
\usepackage{colortbl}
%% Table highlighting: ours (blue, the colour of c = 1.1 in the figures), reference (grey), failure or cost (red).
\definecolor{hlours}{HTML}{DCEBF7}
\definecolor{hlref}{HTML}{ECECEC}
\definecolor{hlbad}{HTML}{F9DEDC}
\newcommand{\hb}{\cellcolor{hlbad}}
%% Lists: labels flush with the paragraph left edge, no indentation.
\setlist[itemize,enumerate]{leftmargin=*,labelindent=0pt,itemindent=0pt,listparindent=0pt,topsep=2pt,partopsep=0pt}
\setlist[description]{leftmargin=0pt,labelindent=0pt,itemindent=0pt}
\graphicspath{{figures/}}

\newcommand{\sys}{\textsc{Ariadne}}
\newcommand{\Adm}{\mathrm{Adm}}
\newcommand{\cser}{$c$-serializable}
\newcommand{\dist}{d}
\newcommand{\Ds}{D_{\mathrm{s}}}
\newcommand{\rhoa}{\rho}
\newcommand{\Rc}{R_c}
\newcommand{\Read}{\textsc{Read}}
\newcommand{\pol}[1]{\textsf{#1}}
\algrenewcommand\algorithmicindent{0.9em}
\emergencystretch=1.5em

%% ---------------------------------------------------------------------
%% Results bridge.  Every measured number in the text is \res[shape]{key}
%% (or \ra, \rb, \ri, \rarng).  paper/make_results.py computes the keys
%% from the result files and writes results.tex (\resdef{key}{value}).
%% A key without a value prints its shape in orange in a draft; with
%% \resultsfinaltrue it is a compile error, so no missing number can
%% survive into a final build.
%% ---------------------------------------------------------------------
\newif\ifresultsfinal
\resultsfinalfalse
\newwrite\reskeys
\immediate\openout\reskeys=\jobname.reskeys
\newcommand{\resdef}[2]{\expandafter\gdef\csname res/#1\endcsname{#2}}
\DeclareRobustCommand{\res}[2][XX]{%
  \immediate\write\reskeys{#2}%
  \ifcsname res/#2\endcsname\csname res/#2\endcsname\else
    \ifresultsfinal\PackageError{ariadne}{No result for key #2}{Run paper/make_results.py on the results.}\fi
    {\color{orange}#1}%
  \fi}
%% A value computed from smoke-scale results (make_results.py --suffix _smoke): orange in a draft, an error in a final build.
\newcommand{\resprov}[1]{\ifresultsfinal\PackageError{ariadne}{Smoke-scale value in a final build}{Run paper/make_results.py on the full results.}\fi{\color{orange}#1}}
\newcommand{\ra}[1]{\res[00.0]{#1}}   % one decimal
\newcommand{\rb}[1]{\res[0.00]{#1}}   % two decimals
\newcommand{\ri}[1]{\res[00]{#1}}     % integer
%% A range of two keys at one decimal; prints one value when both round to the same printed number.
\newcommand{\rarng}[2]{\immediate\write\reskeys{#2}%
  \ifcsname res/#1\endcsname\ifcsname res/#2\endcsname
    \expandafter\let\expandafter\rarngA\csname res/#1\endcsname\expandafter\let\expandafter\rarngB\csname res/#2\endcsname
    \ifx\rarngA\rarngB\ra{#1}\else\ra{#1} to \ra{#2}\fi
  \else\ra{#1} to \ra{#2}\fi\else\ra{#1} to \ra{#2}\fi}
%% Figures are drawn from the results by ariadne.plots (paper/build.sh).  A figure that is not there yet is a framed
%% note in a draft and an error in a final build.
\newcommand{\figstatus}[1]{}
\newcommand{\datafig}[2][\columnwidth]{\IfFileExists{figures/#2.pdf}{\includegraphics[width=#1]{#2}}%
  {\ifresultsfinal\PackageError{ariadne}{Missing figure #2}{Run paper/build.sh on the results.}\fi
   \fbox{\parbox[c][0.45#1][c]{\dimexpr#1-2\fboxsep-2\fboxrule}{\centering\color{orange}\sffamily\small figure \texttt{\detokenize{#2}}: pending}}}}
\InputIfFileExists{results.tex}{}{}

%% The appendix is the full version of the paper (a separate technical report, in the artifact repository).
%% In the 12-page paper the appendix sections are cited as sections of the full version.
\newcommand{\appProofs}{the full version~\cite[App.~A]{ariadnefull}}
\newcommand{\appPG}{the full version~\cite[App.~B]{ariadnefull}}
\newcommand{\appLaw}{the full version~\cite[App.~C]{ariadnefull}}
\newcommand{\appExp}{the full version~\cite[App.~D]{ariadnefull}}
\newcommand{\appAudit}{the full version~\cite[App.~E]{ariadnefull}}
\newcommand{\appRepro}{the full version~\cite[App.~F]{ariadnefull}}
\newcommand{\appRelated}{the full version~\cite[App.~G]{ariadnefull}}
\newcommand{\AppProofs}{The full version~\cite[App.~A]{ariadnefull}}
\newcommand{\AppExp}{The full version~\cite[App.~D]{ariadnefull}}
